\section{An actual finite separation tower for the original NC carrier}
\label{sec:hironaka-normal-crossings-boundary-separation}

The strict carrier in this section is the actual closure of the preimage
off the differential exceptional locus. The construction stops at original
branch count at most one. All centres have original count at least two;
the Cartier centre-removal case is never used in this tower.

\begin{lemma}[Actual original density off a higher highest stratum]
\label{hir:general-nc-higher-highest-relative-density}
\lean{Hironaka.HironakaNormalCrossingsBoundarySeparation.subset_closure_sdiff_highest_of_two_le,Hironaka.HironakaNormalCrossingsBoundarySeparation.geoStrictTransform_highest_eq_strictTransform}
\leanok
\uses{hir:general-nc-highest-stratum-branch-adapted-coordinates,hir:general-nc-intrinsic-branch-multiplicity}
Let $H$ be a genuine original NC carrier whose branch counts are at most
$r\ge2$, and let $C$ be its whole highest stratum. Then
$H\subseteq\overline{H\setminus C}$. Consequently the centre-removed
geometric strict transform in any genuine blowup of $C$ equals its
actual differential-exceptional-locus strict transform.
\end{lemma}
\begin{proof}
A point of $H$ outside $C$ already belongs to $H\setminus C$.
At a point of $C$, the produced actual centred coordinates express
$H$ as the union of the last $r$ coordinate hyperplanes and $C$ as
their intersection. There are two distinct indices $i,j$ in that block
since $r\ge2$. For any neighbourhood of the original point, the actual
coordinate inverse is continuous at zero and has an open target domain.
Hence for sufficiently small nonzero complex $t$, the vector $t e_j$
lies in that target and its actual inverse lies in the chosen
neighbourhood. Its $i$th coordinate is zero, so it belongs to the
actual original carrier; its $j$th coordinate is nonzero, so it lies
outside the whole centre. Every neighbourhood thus meets
$H\setminus C$, proving the actual relative density. The proved
genuine blowup strict-transform comparison applies to this produced
density and gives the stated exact equality. No density or strict
transform dictionary is an input.
\end{proof}

\begin{lemma}[Actual smooth terminal carrier and transverse boundary]
\label{hir:general-nc-count-one-smooth-terminal-pair}
\lean{Hironaka.HironakaNormalCrossingsBoundarySeparation.regularAtDim_of_branchMultiplicity_le_one,Hironaka.HironakaNormalCrossingsBoundarySeparation.isClosedSubmanifold_of_branchMultiplicity_le_one,Hironaka.HironakaNormalCrossingsBoundarySeparation.isTransverse_of_boundaryAdapted_count_le_one,Hironaka.HironakaNormalCrossingsBoundarySeparation.isSNCDivisor_append_original_of_count_le_one}
\leanok
\uses{hir:general-nc-simultaneous-original-boundary,hir:general-nc-highest-multiplicity-smooth-centre}
If a boundary adapted pair has all original branch counts at most one,
its whole original carrier is a closed smooth submanifold, regular
of dimension $n-1$ at each of its points, and is transverse to its
actual SNC boundary. The list obtained by appending that entire
original carrier to the boundary is SNC.
\end{lemma}
\begin{proof}
At a point of the carrier, the actual intrinsic count is positive,
so it is exactly one. The intrinsic count-one superlevel is exactly
the whole carrier, including all its connected components, because
positivity is equivalent to actual carrier membership. The genuine
highest-superlevel regularity theorem therefore produces its actual
regular coordinates of dimension $n-1$ everywhere on the carrier;
it is closed by the original pair hypothesis. Alternatively in the
actual simultaneous coordinates its selected original set has exactly
one index. The whole union equation is then that single hyperplane
equation, hence also its singleton coordinate-subspace equation.
The existing through boundary singletons are disjoint from it and
mutually distinct, giving actual transversality on the same domain.
Appending the original carrier adds precisely that new singleton at
points of the carrier and no through member elsewhere, so the actual
whole appended list satisfies SNC. If the original carrier is empty,
all regularity and transversality claims are vacuous and the appended
empty member creates no through coordinate condition. Dimension zero
is included: positive original count there is impossible.
\end{proof}

\begin{definition}[A genuine intrinsic higher separation step]
\label{hir:general-nc-intrinsic-higher-separation-step}
\lean{Hironaka.HironakaNormalCrossingsBoundarySeparation.IsHigherSeparationCentre}
\leanok
\uses{hir:general-nc-simultaneous-original-boundary}
A current carrier $H$, current finite boundary $E$ and current centre
$C$ form an intrinsic higher separation step when they have an actual
boundary adapted pair, an attained whole global original maximum
$r\ge2$, and $C$ is exactly that whole highest original stratum.
This is a geometric condition on one current centre, not a finite
tower or a termination certificate.
\end{definition}

\begin{lemma}[A genuinely produced higher step and exact original transform]
\label{hir:general-nc-actual-higher-separation-step}
\lean{Hironaka.HironakaNormalCrossingsBoundarySeparation.exists_higher_boundary_blowup,Hironaka.HironakaNormalCrossingsBoundarySeparation.higherSeparationCentre_subset,Hironaka.HironakaNormalCrossingsBoundarySeparation.higherSeparationCentre_geometry}
\leanok
\uses{hir:general-nc-boundary-transform-original-count-step,hir:general-nc-higher-highest-relative-density}
A boundary adapted pair whose original counts are not all at most
one internally produces its attained maximum $r\ge2$, its nonempty
whole closed smooth centre of codimension $r$, a genuine blowup, and
an actual transformed pair with every original count at most $r-1$.
Its original carrier is exactly the differential strict transform;
the current accumulated boundary is exactly the strict old list
followed by the new exceptional preimage.
\end{lemma}
\begin{proof}
The failure of the count-one bound gives an actual point of count
greater than one and hence a nonempty original carrier. The genuine
finite intrinsic-count image produces an attained positive maximum
$r$, bounding that point, so $r\ge2$. Its whole stratum is nonempty,
closed and regular of dimension $n-r$ at every centre point, by the
actual maximum-stratum producer. Its entire centre lies in the
original carrier because its count is positive. The actual simultaneous
pair coordinates flatten the centre using the original selected indices;
all through old boundary singletons are disjoint from those indices,
so the entire centre is transverse to the actual SNC boundary.
The genuine blowup and whole
transformed-pair producer then give the actual projection, true
accumulated boundary list and smaller original count. The preceding
actual relative-density theorem identifies this centre-removed
original strict carrier with the differential strict carrier.
The old boundary uses exactly its genuine centre-removed strict
transforms, as required by the accumulated-boundary operation; no
replacement of those boundary definitions is made.
\end{proof}

\begin{theorem}[A bounded actual tower with literal higher centres]
\label{hir:general-nc-bounded-actual-separation-tower}
\lean{Hironaka.HironakaNormalCrossingsBoundarySeparation.exists_boundary_separation_tower_bound,Hironaka.HironakaNormalCrossingsBoundarySeparation.stageGeometry_of_higherSeparationCentres}
\leanok
\uses{hir:general-nc-actual-higher-separation-step,hir:general-nc-count-one-smooth-terminal-pair,hir:blowup-tower}
For a boundary adapted pair with all original counts bounded by $b$,
there is an actual finite blowup tower of length at most
$\max(b-1,0)$ whose whole differential strict original carrier has
count at most one and is boundary adapted with its literal accumulated
boundary. Every actual recorded centre is an intrinsic higher
separation step for the actual current differential strict carrier
and actual current accumulated boundary. In particular the boundary
at each actual centre stage is SNC and every centre has codimension
at least two and lies in the current strict original carrier.
\end{theorem}
\begin{proof}
Induct on the actual integer bound $b$, uniformly in the current
manifold, original carrier and current finite boundary. If all actual
counts are at most one, take the real empty tower. Its projection is
the identity, its strict original carrier is the same closed carrier,
and its boundary list is unchanged. It has length zero and no centre
conditions. This covers an empty carrier and the base case $b=0$.

Otherwise apply the preceding produced higher step. Its attained
maximum satisfies $2\le r\le b$. The new actual original count is
at most $r-1\le b-1$. Apply the induction hypothesis to the actual
blowup manifold and produced transformed pair, then prepend the
actual first blowup to the resulting real tail tower. The exact
first-step differential strict-transform dictionary and the genuine
holomorphic strict-composition theorem identify every tail current
carrier with the actual differential strict transform under its
whole prefix. The actual accumulated-boundary append identity
identifies every tail boundary with the literal original boundary
transformed along that whole prefix. Thus all actual tail centre
conditions transport to the concatenated tower; the first centre
has the produced intrinsic higher-step property. Tower lengths add,
and the tail bound gives at most $1+\max(b-2,0)=b-1$, since this
case has $b\ge2$. The final actual pair and count-one bound are
transported by those same exact dictionaries. This is an actual
finite construction, with no caller-supplied candidate tower,
termination certificate or global ordering of original branches.
\end{proof}

\begin{theorem}[A genuine global NC separation and complement isomorphism]
\label{hir:general-nc-global-separation-complement-isomorphism}
\lean{Hironaka.HironakaNormalCrossingsBoundarySeparation.centresOver_of_higherSeparationCentres,Hironaka.HironakaNormalCrossingsBoundarySeparation.isIsomorphismOver_compl_of_higherSeparationCentres,Hironaka.HironakaNormalCrossingsBoundarySeparation.exists_boundary_separation_tower,Hironaka.HironakaNormalCrossingsBoundarySeparation.exists_length_zero_separation_of_empty,Hironaka.HironakaNormalCrossingsBoundarySeparation.exists_normalCrossings_separation_tower}
\leanok
\uses{hir:general-nc-bounded-actual-separation-tower,hir:general-nc-boundary-empty-producer,hir:general-nc-count-one-smooth-terminal-pair}
Every actual original NC carrier has a genuine separation tower of
length at most $\max(n-1,0)$ starting with empty boundary. Its final
whole differential strict carrier is smooth, transverse to the
actual accumulated SNC boundary, and appending that carrier gives
a genuine SNC list. More generally an initial boundary adapted pair
has the same conclusion with its true accumulated old boundary.
The tower projection is an isomorphism over the complement of the
original carrier.
\end{theorem}
\begin{proof}
The genuine actual count is always at most $n$, so apply the bounded
tower producer with $b=n$. Its actual terminal pair and count-one
bound give smoothness, pointwise dimension, transversality and the
SNC appended list by the terminal-pair theorem. For an original NC
carrier alone, its empty-boundary pair is internally produced first.

Each literal centre lies in its actual current differential strict
carrier. Closedness of the original carrier and continuity of the
actual prefix projection put that entire strict carrier inside the
preimage of the original carrier. Thus every real recorded centre
lies over the original carrier. Restricting the actual tower to its
open complement makes every centre empty; the resulting true tower
has no differential exceptional point there. The actual exceptional
image therefore lies in the original carrier. The genuine blowup
sequence is surjective and injective away from its exceptional locus,
so the proved blowup-sequence isomorphism theorem gives the stated
isomorphism over that whole open complement. No compactness,
connectedness or finite global component family is required.
\end{proof}

\begin{remark}
The output is an actual intrinsic geometric separation tower, with
literal current carriers and accumulated boundary at every stage.
This does not imply cosupport admissibility for a separately marked
ideal. The old ordered invariant retains its colexicographic label
tie-break, and its canonical marked-ideal sequence is not identified
with this geometric tower. That full marked-ideal general-NC descent
remains open.
\end{remark}
